% Generated draft: requires mathematical and source review.
\documentclass[11pt]{article}
\usepackage[margin=25mm]{geometry}
\usepackage{fontspec}
\setmainfont{Latin Modern Roman}
\usepackage{amsmath,amssymb,mathtools,unicode-math}
\setmathfont{Latin Modern Math}
\usepackage{enumitem,needspace,xurl,hyperref}
\setlength{\parindent}{0pt}
\setlength{\parskip}{5pt}
\setlength{\emergencystretch}{3em}
\newcommand{\N}{\mathbb{N}}
\newcommand{\Z}{\mathbb{Z}}
\newcommand{\Q}{\mathbb{Q}}
\newcommand{\R}{\mathbb{R}}
\newcommand{\C}{\mathbb{C}}
\newcommand{\F}{\mathbb{F}}
\newcommand{\A}{\mathbb{A}}
\newcommand{\PP}{\mathbb{P}}
\newcommand{\E}{\mathbb{E}}
\newcommand{\eps}{\varepsilon}
\newcommand{\dd}{\,\mathrm{d}}
\begin{document}
\section*{UnsolvedMath ID: 30005772}
\textbf{Source problem code:} catalogue entry
\subsection{Definitions and mathematical statement}
For an integer \(n\geq 0\), let \([n]=\{1,\ldots,n\}\), with \([0]=\varnothing\), and let \(S_n\) be the set of permutations of \([n]\). For \(\pi\in S_n\), write
\[
\operatorname{Fix}(\pi)=\{i\in[n]:\pi(i)=i\},\qquad f(\pi)=|\operatorname{Fix}(\pi)|.
\]
For an integer \(k\geq0\), a permutation with colored fixed points is a pair \((\pi,c)\), where \(\pi\in S_n\) and \(c:\operatorname{Fix}(\pi)\to[k]\); nonfixed points receive no colors. Define the polynomial
\[
A(n,t)=\sum_{\pi\in S_n}t^{f(\pi)}.
\]
Here \(t^0=1\), including at \(t=0\). Thus \(A(n,k)\) counts these pairs for nonnegative integers \(k\), and the same polynomial defines \(A(n,k)\) for every integer \(k\). We use the empty-permutation convention \(A(0,t)=1\); the supplied record does not explicitly specify the boundary convention at \(n=0\).

The supplied record poses the following questions.
\begin{enumerate}
\item For every integer \(k<0\), find a combinatorial interpretation of \(A(n,k)\), valid for all positive integers \(n\).
\item For every integer \(k\in\Z\), find a combinatorial interpretation of \(A(2n,k)\), valid for all positive integers \(n\).
\end{enumerate}
The source presents the second question as potentially easier. It additionally asserts
\[
A(2n,t)\geq0\qquad(n\geq1,\ t\in\R),
\]
with justification attributed to a moment-integral representation. This is a supplied source assertion; the extracted integral and its explanatory exponent are corrupted, so an exact integral formula is not adopted here.

The phrase ``combinatorial interpretation'' is not defined in the supplied record. In particular, the first question does not specify whether signed or weighted enumeration, a sign normalization, or another interpretation is permitted. Accordingly, these are interpretation requests with unresolved admissibility criteria, rather than fully specified assertions that the values are cardinalities of finite sets.
\subsection{Short English statement}
The polynomial \(A(n,k)\) counts permutations whose fixed points independently receive one of \(k\) colors when \(k\geq0\). The questions ask for combinatorial meanings at negative integer parameters, and at all integer parameters when the permutation size is even. The permitted meaning of a combinatorial interpretation, especially any use of signs or weights, requires clarification.
\subsection{Sources}
\begin{itemize}
\item Supplied local UnsolvedMath registry record, permanent ID \(30005772\), title ``Combinatorial Interpretations of Negative k-Arrangements,'' exact source problem code \texttt{catalogue entry}, in the Oberwolfach Reports Open Problems collection. The record gives proposed year 2024 and does not name a proposer.
\item \emph{Mini-Workshop: Permutation Patterns} (2024), Oberwolfach Reports, pp. 273--308, DOI \texttt{10.4171/owr/2024/6}. Supplied source citation and literature-source label: ``original OWR report.'' \url{https://doi.org/10.4171/owr/2024/6}.
\item Blitvić--Steingrímsson, \emph{Permutations, moments, measures}. Supplied literature source: \url{https://arxiv.org/abs/2001.00280}. The supplied record does not identify a specific passage supporting these questions.
\item Supplied literature assessment dated \texttt{2026-08-22}: the report is said to pose both interpretation questions explicitly, and the cited literature is said to record them as unresolved. This is supplied dated curation, not an independently verified literature assessment. The background expressly requests independent verification of time-sensitive status and cited claims.
\item The registry labels the problem ``open'' and the statement \texttt{corrected\_verified}; its verification note reports that the official report was checked directly and that definitions were restored while extraction artifacts were removed. These are supplied provenance claims, not verification performed for this draft.
\end{itemize}
\end{document}
